% math0 manuscript. Build with: latexmk -pdf paper.tex
\documentclass[11pt]{amsart}
\usepackage[margin=1.1in]{geometry}
\usepackage{amsmath,amssymb,amsthm,mathtools}
\usepackage{url}
\usepackage{hyperref}
\usepackage{xcolor}
\usepackage{listings}
\usepackage{newunicodechar}

\lstdefinelanguage{lean}{
  keywords={theorem,lemma,def,abbrev,example,fun,by,have,show,from,let,in,if,then,else,match,with,
    namespace,end,open,import,section,variable,instance,structure,where,calc,obtain,intro,exact,
    refine,rcases,rintro,simp,omega,decide,norm_num,linarith,nlinarith,positivity,induction,cases,
    private,rw,ring,linear_combination,by_contra},
  sensitive=true, comment=[l]{--}, morecomment=[s]{/-}{-/}, string=[b]"}
\lstset{language=lean, basicstyle=\ttfamily\footnotesize, columns=fullflexible, keepspaces=true,
  breaklines=true, keywordstyle=\color{blue!60!black}\bfseries, commentstyle=\color{gray},
  frame=single, extendedchars=true, inputencoding=utf8,
  literate={ℕ}{{$\mathbb{N}$}}1 {ℤ}{{$\mathbb{Z}$}}1 {∀}{{$\forall$}}1 {∃}{{$\exists$}}1
    {→}{{$\to$}}1 {←}{{$\leftarrow$}}1 {≠}{{$\ne$}}1 {∧}{{$\wedge$}}1 {∣}{{$\mid$}}1
    {≤}{{$\le$}}1 {⟨}{{$\langle$}}1 {⟩}{{$\rangle$}}1 {₀}{{$_0$}}1 {·}{{$\cdot$}}1}
% Unicode commonly found in Lean sources.
\newunicodechar{ℕ}{\ensuremath{\mathbb{N}}}
\newunicodechar{ℤ}{\ensuremath{\mathbb{Z}}}
\newunicodechar{ℚ}{\ensuremath{\mathbb{Q}}}
\newunicodechar{ℝ}{\ensuremath{\mathbb{R}}}
\newunicodechar{ℂ}{\ensuremath{\mathbb{C}}}
\newunicodechar{∀}{\ensuremath{\forall}}
\newunicodechar{∃}{\ensuremath{\exists}}
\newunicodechar{→}{\ensuremath{\to}}
\newunicodechar{←}{\ensuremath{\leftarrow}}
\newunicodechar{↔}{\ensuremath{\leftrightarrow}}
\newunicodechar{↦}{\ensuremath{\mapsto}}
\newunicodechar{¬}{\ensuremath{\neg}}
\newunicodechar{∧}{\ensuremath{\wedge}}
\newunicodechar{∨}{\ensuremath{\vee}}
\newunicodechar{≤}{\ensuremath{\le}}
\newunicodechar{≥}{\ensuremath{\ge}}
\newunicodechar{≠}{\ensuremath{\ne}}
\newunicodechar{∈}{\ensuremath{\in}}
\newunicodechar{∉}{\ensuremath{\notin}}
\newunicodechar{⊆}{\ensuremath{\subseteq}}
\newunicodechar{∩}{\ensuremath{\cap}}
\newunicodechar{∪}{\ensuremath{\cup}}
\newunicodechar{∑}{\ensuremath{\sum}}
\newunicodechar{∏}{\ensuremath{\prod}}
\newunicodechar{⟨}{\ensuremath{\langle}}
\newunicodechar{⟩}{\ensuremath{\rangle}}
\newunicodechar{·}{\ensuremath{\cdot}}
\newunicodechar{∘}{\ensuremath{\circ}}
\newunicodechar{⁻}{\ensuremath{^{-}}}
\newunicodechar{¹}{\ensuremath{^{1}}}
\newunicodechar{²}{\ensuremath{^{2}}}
\newunicodechar{₀}{\ensuremath{_{0}}}
\newunicodechar{₁}{\ensuremath{_{1}}}
\newunicodechar{₂}{\ensuremath{_{2}}}
\newunicodechar{α}{\ensuremath{\alpha}}
\newunicodechar{β}{\ensuremath{\beta}}
\newunicodechar{ε}{\ensuremath{\varepsilon}}
\newunicodechar{δ}{\ensuremath{\delta}}
\newunicodechar{π}{\ensuremath{\pi}}
\newunicodechar{σ}{\ensuremath{\sigma}}
\newunicodechar{φ}{\ensuremath{\varphi}}
\newunicodechar{∞}{\ensuremath{\infty}}
\newunicodechar{⌊}{\ensuremath{\lfloor}}
\newunicodechar{⌋}{\ensuremath{\rfloor}}
\newunicodechar{⌈}{\ensuremath{\lceil}}
\newunicodechar{⌉}{\ensuremath{\rceil}}
\newunicodechar{∣}{\ensuremath{\mid}}
\newunicodechar{≡}{\ensuremath{\equiv}}
\newunicodechar{×}{\ensuremath{\times}}
\newunicodechar{⧸}{\ensuremath{/}}
\newunicodechar{‖}{\ensuremath{\|}}
\newunicodechar{√}{\ensuremath{\surd}}
\newunicodechar{∅}{\ensuremath{\emptyset}}
\newunicodechar{⦃}{\ensuremath{\{\!|}}
\newunicodechar{⦄}{\ensuremath{|\!\}}}
\newunicodechar{𝓝}{\ensuremath{\mathcal{N}}}
\newunicodechar{≫}{\ensuremath{\gg}}
\newunicodechar{≪}{\ensuremath{\ll}}

\theoremstyle{plain}
\newtheorem{theorem}{Theorem}
\newtheorem{lemma}[theorem]{Lemma}
\newtheorem{proposition}[theorem]{Proposition}
\newtheorem{corollary}[theorem]{Corollary}
\newtheorem{conjecture}[theorem]{Conjecture}
\theoremstyle{definition}
\newtheorem{definition}[theorem]{Definition}
\newtheorem{remark}[theorem]{Remark}

\emergencystretch=3em
% Lean identifiers in running text (breakable at dots and underscores).
\expandafter\def\expandafter\UrlBreaks\expandafter{\UrlBreaks\do\_\do\.}
\DeclareUrlCommand\decl{\urlstyle{tt}}
\newcommand{\N}{\mathbb{N}}
\newcommand{\Z}{\mathbb{Z}}
\newcommand{\Q}{\mathbb{Q}}

\title[Heronian triangles with integer circumradius]{No primitive Heronian triangle has an integer circumradius (OEIS A392943)}
\author{math0}
\address{Produced by Claude (an AI model developed by Anthropic) in the math0 project}
\date{October 19, 2026}

\begin{document}
\begin{abstract}
We prove that no primitive Heronian triangle (integer sides with greatest common divisor $1$ and integer area) has an integer circumradius, which is the conjecture recorded by Felix Huber in OEIS A392943 in February 2026; more precisely, every triangle with integer sides and integer circumradius has three even sides.
The result is not new: it follows from R.~van Luijk's 2008 characterization of the circumdiameters of Heron triangles, whose proof states explicitly that a Heron triangle with integer circumradius has even sides. Our proof is a short self-contained argument modulo $4$, and we classify this note as the first formal proof of a known result.
The conjecture, in the form with integer area as in the OEIS entry, is verified in Lean~4 against a statement that math0 formalized from the entry, reviewed for fidelity and frozen before the proof attempt; the definitions of the entry's sequence built on the same notions are checked by the Lean kernel against the first $30$ listed terms.
\end{abstract}
\maketitle

\section{Introduction}

A \emph{Heronian} (or Heron) triangle is a triangle whose side lengths $a,b,c$ and whose area $K$ are integers; it is \emph{primitive} if $\gcd(a,b,c)=1$. Heron's formula
\begin{equation}\label{eq:heron}
  16K^2=P(a,b,c)\coloneqq(a+b+c)(-a+b+c)(a-b+c)(a+b-c)
\end{equation}
expresses the area through the sides, and the circumradius is
\begin{equation}\label{eq:circ}
  R=\frac{abc}{4K}.
\end{equation}
For example, the triangle $(3,4,5)$ is Heronian with $K=6$ and $R=5/2$, while its double $(6,8,10)$ has $K=24$ and $R=5$.

The OEIS sequence A392942 \cite{A392942} lists the semiperimeters of the triangles with integer sides, integer area and integer circumradius ($12,24,30,36,40,\dots$), and A392943 \cite{A392943} counts, for the $n$-th such semiperimeter, the triangles (up to congruence) with that semiperimeter. Both entries were contributed by Felix Huber on February~4, 2026. The entry A392943 contains the comment ``Conjecture: There is no primitive Heronian triangle with integer circumradius.'' All triangles given as examples in the two entries, namely $(30,30,48)$, $(24,70,74)$, $(28,60,80)$, $(42,56,70)$ and $(14,30,40)$, have three even sides. The entries contain no proof and no reference for the conjecture.

\subsection*{Prior work}
The conjecture is a consequence of published results. A.-V.~Kramer and F.~Luca \cite[Proposition~6]{KL} proved that the set of positive integers $k$ that are the circumradius of a Heron triangle has asymptotic density $1$, and that no Heron triangle has a circumradius that is a power of $2$ or a power of a prime $p\equiv 11\pmod{12}$. The case of primes $p\equiv11\pmod{12}$ was also stated as a conjecture by M.~Aassila \cite{Aa}, as recorded in \cite{vL}; Aassila's paper appeared in 2001, after \cite{KL}, and was later marked as retracted by Elemente der Mathematik, whose editorial board noted its strong similarity to \cite{KL}.

R.~van Luijk \cite{vL} proved that a positive integer $k$ is the circumradius of a Heron triangle if and only if it is the circumdiameter of a Heron triangle, if and only if $k$ has a prime divisor $p\equiv1\pmod 4$ \cite[Theorem~1]{vL}. The key step is the implication (4)$\Rightarrow$(5) of his Theorem~2: if a triangle has integer sides $a,b,c$ and integer circumdiameter $d$ with $\gcd(a,b,c,d)=1$, then $d>1$ and every prime divisor of $d$ is $\equiv1\pmod4$. Its Corollary~5 states that if an integer $d$ is the circumdiameter of a triangle with integer sides $a,b,c$ and $p\not\equiv1\pmod 4$ is prime, then $p^{v_p(d)}$ divides $\gcd(a,b,c)$. The proof of the implication (1)$\Rightarrow$(2) of Theorem~1 in \cite{vL} begins: ``let $a$, $b$, and $c$ be the sides of a Heron triangle with circumradius $k$. Then the circumdiameter $2k$ is even and by Corollary~5, so are $a$, $b$, and $c$.'' This is exactly the statement that no primitive Heronian triangle has an integer circumradius. Van Luijk's argument rests on the identity $(4K)^2+(a^2+b^2-c^2)^2=4a^2b^2$ and on the fact that $p^{2m}\mid s^2+t^2$ forces $p^m\mid s$ and $p^m\mid t$ for primes $p\not\equiv1\pmod4$ \cite[Lemma~4]{vL}. He also shows that a triangle with integer sides and rational area is Heronian \cite[Lemma~3]{vL}; by \eqref{eq:circ}, a triangle with integer sides and integer circumradius therefore has integer area automatically. E.~Read \cite{Read} has more recently studied integer triangles with integer circumradii; the publicly available extract of that article notes that integer area is necessary for an integer circumradius, and we have not consulted its full text.

Neither OEIS entry cites these works, the conjecture is not among the OEIS statements of the formal-conjectures repository \cite{FC} (checked at commit \texttt{c77fdd4}), and we are not aware of an earlier formal proof.

\subsection*{This note}
We give a self-contained proof that integer circumradius forces all sides to be even (Lemma~\ref{lem:key} and Remark~\ref{rem:general}). It is a weak form of the case $p=2$ of van Luijk's Corollary~5 (we obtain $2\mid\gcd(a,b,c)$ rather than $2^{v_2(2R)}\mid\gcd(a,b,c)$), proved directly from Heron's formula by an argument modulo $4$. We classify the result as \emph{a known result, formalized here for the first time}: the mathematics has been available since 2008, and the contribution of this note is a written resolution of the OEIS conjecture that is linked to the literature, together with a kernel-checked proof of a frozen formal statement.

\section{Statement}\label{sec:statement}

The name of A392943 \cite{A392943} is, verbatim, ``a(n) is the number of integer-sided triangles whose area and circumradius are integers, and whose semiperimeter equals A392942(n).'' The entry is authored by Felix Huber and dated February~4, 2026, and its COMMENTS section contains, verbatim, the following comment:
\begin{quote}\raggedright
\texttt{Conjecture: There is no primitive Heronian triangle with integer circumradius.}
\end{quote}

\begin{conjecture}[Huber \cite{A392943}]\label{conj}
If a triangle has integer sides $a,b,c$ with $\gcd(a,b,c)=1$ and integer area, then its circumradius is not an integer.
\end{conjecture}

\begin{theorem}\label{thm:main}
Let $x,y,z$ be positive integers with $x<y+z$, $y<x+z$, $z<x+y$ and $\gcd(x,y,z)=1$, and let $K\in\N$ satisfy $16K^2=P(x,y,z)$, where $\N=\{0,1,2,\dots\}$. Then there is no $R\in\N$ with $4KR=xyz$. In particular, Conjecture~\ref{conj} holds.
\end{theorem}

Theorem~\ref{thm:main} is the formal statement below, written in standard notation. The gcd hypothesis cannot be dropped: $(6,8,10)$ with $K=24$ and $R=5$ satisfies all other hypotheses and $4KR=xyz$.

\subsection*{Lean formalization}
The statement was formalized by math0 in \path{lean/Math0/Statements/OEIS/A392943.lean}, reviewed for fidelity and frozen (SHA-256 \texttt{efec1ccc\ldots 893a3a4a}) before any proof attempt. The definitions it uses and the statement are, verbatim except that docstrings are omitted:
\begin{lstlisting}
namespace Math0.Statements.A392943

def IsTriangle (x y z : ℕ) : Prop :=
  0 < x ∧ 0 < y ∧ 0 < z ∧ x < y + z ∧ y < x + z ∧ z < x + y

def heron16 (x y z : ℕ) : ℤ :=
  ((x : ℤ) + y + z) * (-x + y + z) * (x - y + z) * (x + y - z)

def HasArea (x y z K : ℕ) : Prop :=
  16 * (K : ℤ) ^ 2 = heron16 x y z

def statement : Prop :=
  ∀ x y z K R : ℕ, IsTriangle x y z → HasArea x y z K → Nat.gcd x (Nat.gcd y z) = 1 →
    4 * K * R ≠ x * y * z

end Math0.Statements.A392943
\end{lstlisting}
\decl{IsTriangle} says that $x,y,z$ are the sides of a nondegenerate triangle. \decl{heron16} is $P(x,y,z)$, computed in $\Z$: the coercion of $x$ to $\Z$ forces every leaf of the expression to be coerced, so there is no truncated subtraction of natural numbers. \texttt{HasArea x y z K} says $16K^2=P(x,y,z)$, i.e.\ by \eqref{eq:heron} that the area is the natural number $K$; for a nondegenerate triangle $P(x,y,z)>0$, so $K>0$. Since $K>0$, the circumradius $xyz/(4K)$ is an integer if and only if $4KR=xyz$ for some integer $R$, and such an $R$ is positive, so quantifying over $R\in\N$ loses nothing. Primitivity is $\gcd(x,\gcd(y,z))=1$. No hypotheses are added, and the hypotheses are satisfiable (by $(3,4,5)$ with $K=6$), so the statement is not vacuous.

The statement file also defines the two sequences from the same notions and checks them against the OEIS data. The predicate \texttt{IsGood x y z} is \texttt{IsTriangle x y z} together with the existence of $K,R\in\N$ with \texttt{HasArea x y z K} and $4KR=xyz$, matching the name of the entry and the Maple test of the entry (\texttt{issqr(f)} and $xyz\equiv0\pmod{4\sqrt f}$, with $f=P/16$). The function \texttt{count s} counts the triangles $x\le y\le z$ with $x+y+z=2s$ satisfying \decl{IsGood}, using the loop bounds of the Maple program, and \decl{count_eq_card} proves that it equals the cardinality of the corresponding finite set of triples. The lemma \decl{even_perimeter} shows that a triangle with integer sides and integer area has even perimeter, so no triangle is lost by taking the semiperimeter $s$ to be a natural number. The term \texttt{A392942 n} is computed by a search over $s\le 12n$ (every positive multiple of $12$ is a semiperimeter, by the triangles $(6m,8m,10m)$), and \decl{A392942_eq_nth} proves that for $n\ge1$ it is the $n$-th smallest semiperimeter, \texttt{Nat.nth IsSemiperimeter (n - 1)}, matching offset $1$. Finally
\begin{lstlisting}
def a (n : ℕ) : ℕ := count (A392942 n)

theorem a_data : (List.range 30).map (fun i => a (i + 1)) =
    [1, 1, 1, 1, 1, 1, 1, 1, 1, 2, 1, 1, 1, 2, 3, 2, 1, 1, 2, 1, 1, 4, 2, 2, 3, 2, 3, 1, 2,
      2] := by
  decide +kernel
\end{lstlisting}
so the Lean kernel itself evaluates $a(1),\dots,a(30)$ and compares them with the first $30$ of the $88$ listed terms of A392943 (these involve the semiperimeters up to $A392942(30)=156$). For speed, the kernel evaluates a Boolean test \decl{goodB}, which combines a quadratic-residue sieve with \decl{Nat.sqrt}; the file proves \decl{goodB_iff}, i.e.\ that \texttt{goodB x y z = true} if and only if \texttt{IsGood x y z}. An independent Python computation by the math0 reviewers, which is not part of the formal verification, reproduced all $88$ listed terms of A392943 and all $56$ listed terms of A392942 from the same definition. The sequence definitions do not enter \decl{statement}; they tie the notions \decl{IsTriangle}, \decl{heron16} and \decl{HasArea} used by \decl{statement} to the OEIS data.

\section{Proof}\label{sec:proof}

The proof follows the Lean file \path{lean/Math0/Frontier/OEIS/A392943.lean}. It consists of one lemma, \decl{two_dvd_of_heron} (Lemma~\ref{lem:key}), and the main theorem \decl{result}. We write $P(a,b,c)$ for the polynomial in \eqref{eq:heron}, now with $a,b,c$ in $\Z$.

\begin{lemma}\label{lem:identity}
For all integers $a,b,c$, with $D=b^2+c^2-a^2$,
\[
  P(a,b,c)=4b^2c^2-D^2 .
\]
Moreover $P$ is symmetric: $P(a,b,c)=P(b,a,c)=P(c,a,b)$.
\end{lemma}

\begin{proof}
Grouping the factors, $P(a,b,c)=\bigl((b+c)^2-a^2\bigr)\bigl(a^2-(b-c)^2\bigr)=\bigl(2bc+D\bigr)\bigl(2bc-D\bigr)$. The second claim holds because the four factors of $P(b,a,c)$ and of $P(c,a,b)$ are a permutation of the four factors of $P(a,b,c)$.
\end{proof}

\begin{lemma}\label{lem:sq}
Let $u,v\in\Z$. If $u^2\mid v^2$, then $u\mid v$.
\end{lemma}

\begin{proof}
If $u=0$, then $v^2=0$ and $v=0$. Otherwise let $g=\gcd(u,v)>0$ and write $u=gu'$, $v=gv'$ with $\gcd(u',v')=1$. Then $u'^2\mid v'^2$, and $u'^2$ is coprime to $v'^2$, so $u'^2=1$ and $u=\pm g$ divides $v$.
\end{proof}

\begin{lemma}\label{lem:key}
Let $a,b,c,R\in\Z$ with $b\ne0$, $c\ne0$ and
\begin{equation}\label{eq:hyp}
  P(a,b,c)\,R^2=(abc)^2 .
\end{equation}
Then $a$ is even.
\end{lemma}

\begin{proof}
Let $D=b^2+c^2-a^2$. By Lemma~\ref{lem:identity} and \eqref{eq:hyp},
\[
  (bc)^2(4R^2-a^2)-(RD)^2=R^2\bigl(4b^2c^2-D^2\bigr)-(abc)^2=R^2P(a,b,c)-(abc)^2=0,
\]
so
\begin{equation}\label{eq:h1}
  (bc)^2(4R^2-a^2)=(RD)^2 .
\end{equation}
Hence $(bc)^2\mid(RD)^2$, and by Lemma~\ref{lem:sq} there is $m\in\Z$ with $RD=bc\,m$. Substituting into \eqref{eq:h1} gives $(bc)^2(4R^2-a^2)=(bc)^2m^2$. Since $bc\ne0$, we may cancel $(bc)^2$ and obtain
\begin{equation}\label{eq:h3}
  a^2+m^2=4R^2 .
\end{equation}
Suppose that $a$ is odd. Then $a^2\equiv1\pmod4$. If $m=2j$ is even, then $m^2=4j^2$ and $a^2+m^2\equiv1\pmod4$; if $m$ is odd, then $m^2\equiv1\pmod 4$ and $a^2+m^2\equiv2\pmod4$. Both contradict \eqref{eq:h3}, whose right-hand side is divisible by $4$. Hence $a$ is even.
\end{proof}

\begin{proof}[Proof of Theorem~\ref{thm:main}]
Let $x,y,z,K,R\in\N$ satisfy the hypotheses, and suppose $4KR=xyz$. Of the hypothesis \decl{IsTriangle} only $x,y,z>0$ is used. In $\Z$ we have
\[
  P(x,y,z)\,R^2=16K^2R^2=(4KR)^2=(xyz)^2 .
\]
The integers $x,y,z$ are nonzero, so Lemma~\ref{lem:key} applies to $(a,b,c)=(x,y,z)$ and gives $2\mid x$. By the symmetry in Lemma~\ref{lem:identity}, $P(y,x,z)R^2=(yxz)^2$ and $P(z,x,y)R^2=(zxy)^2$, so Lemma~\ref{lem:key} applied to $(y,x,z)$ and to $(z,x,y)$ gives $2\mid y$ and $2\mid z$. Therefore $2$ divides $\gcd(x,\gcd(y,z))=1$, a contradiction.
\end{proof}

\begin{remark}\label{rem:general}
The proof does not use the triangle inequalities, and it uses the integrality of the area only through the identity $P(x,y,z)R^2=(xyz)^2$. For any triangle with integer sides $a,b,c$ and integer circumradius $R$, the (real) area $K>0$ satisfies \eqref{eq:heron} and \eqref{eq:circ}, so $P(a,b,c)R^2=16K^2R^2=(abc)^2$, and Lemma~\ref{lem:key}, applied three times as above, shows that $a$, $b$ and $c$ are all even. Lemma~\ref{lem:key} is formalized in this generality (for arbitrary integers $a,b,c,R$ with $b,c\neq 0$); only the final theorem is stated with $K\in\N$, as in the OEIS conjecture. Geometrically, if $A$ is the angle opposite the side $a$, then $a=2R\sin A$ by the law of sines and $D=2bc\cos A$ by the law of cosines, so $m=RD/(bc)=2R\cos A$ and \eqref{eq:h3} is $\sin^2A+\cos^2A=1$ multiplied by $4R^2$. Van Luijk's Corollary~5 \cite{vL} gives the stronger conclusion that $2^{v_2(2R)}$ divides $\gcd(a,b,c)$.
\end{remark}

\section{Machine verification}\label{sec:lean}

\subsection*{Files and declarations}
The proof is in \path{lean/Math0/Frontier/OEIS/A392943.lean} in the math0 repository, \url{https://github.com/Foundation-0/math0}. It imports the frozen statement file \path{lean/Math0/Statements/OEIS/A392943.lean} (Section~\ref{sec:statement}) and the audit commands in \path{lean/Math0/Audit.lean}. The main declaration is \decl{Math0.Frontier.A392943.result}, whose type is exactly \decl{Math0.Statements.A392943.statement}. The only auxiliary declaration is Lemma~\ref{lem:key}; its statement and proof are, verbatim:
\begin{lstlisting}
private theorem two_dvd_of_heron (a b c R : ℤ) (hb : b ≠ 0) (hc : c ≠ 0)
    (h : (a + b + c) * (-a + b + c) * (a - b + c) * (a + b - c) * R ^ 2 = (a * b * c) ^ 2) :
    2 ∣ a := by
  have h1 : (b * c) ^ 2 * (4 * R ^ 2 - a ^ 2) = (R * (b ^ 2 + c ^ 2 - a ^ 2)) ^ 2 := by
    linear_combination h
  have hdvd : (b * c) ^ 2 ∣ (R * (b ^ 2 + c ^ 2 - a ^ 2)) ^ 2 := ⟨_, h1.symm⟩
  obtain ⟨m, hm⟩ := (Int.pow_dvd_pow_iff (by norm_num)).1 hdvd
  have hbc : (b * c) ^ 2 ≠ 0 := pow_ne_zero _ (mul_ne_zero hb hc)
  have h2 : (b * c) ^ 2 * (4 * R ^ 2 - a ^ 2) = (b * c) ^ 2 * m ^ 2 := by
    rw [h1, hm]; ring
  have h3 : a ^ 2 + m ^ 2 = 4 * R ^ 2 := by
    have := mul_left_cancel₀ hbc h2
    linear_combination -this
  by_contra hodd
  have ha : Odd a := Int.not_even_iff_odd.1 (fun he => hodd (even_iff_two_dvd.1 he))
  have ha2 : a ^ 2 % 4 = 1 := Int.sq_mod_four_eq_one_of_odd ha
  have h4 : (a ^ 2 + m ^ 2) % 4 = 0 := by rw [h3]; simp
  rcases Int.even_or_odd m with ⟨j, rfl⟩ | hm'
  · have hm2 : (j + j) ^ 2 = 4 * (j ^ 2) := by ring
    rw [hm2] at h4
    omega
  · have hm2 : m ^ 2 % 4 = 1 := Int.sq_mod_four_eq_one_of_odd hm'
    omega
\end{lstlisting}
The steps correspond to the proof of Lemma~\ref{lem:key}: \texttt{h1} is \eqref{eq:h1}, obtained by \texttt{linear\_combination}, which checks with \texttt{ring} that the difference of the two sides of \texttt{h1} equals the difference of the two sides of \texttt{h} (this is Lemma~\ref{lem:identity}); Lemma~\ref{lem:sq} is the core Lean lemma \decl{Int.pow_dvd_pow_iff}; \texttt{h3} is \eqref{eq:h3}, obtained by cancelling with \texttt{mul\_left\_cancel}$_0$; and the final case analysis modulo $4$ uses Mathlib's \decl{Int.sq_mod_four_eq_one_of_odd} and the decision procedure \texttt{omega}. The proof of \decl{result} follows the proof of Theorem~\ref{thm:main}: it casts $4KR=xyz$ to $\Z$, derives $P(x,y,z)R^2=(xyz)^2$ by rewriting with \decl{HasArea} and \texttt{ring}, obtains the two permuted hypotheses by \texttt{linear\_combination}, applies \decl{two_dvd_of_heron} three times, casts the three divisibilities back to $\N$, and concludes with \decl{Nat.dvd_gcd} and \texttt{omega}. The proof performs no large computation; \texttt{native\_decide}, \texttt{sorry} and \texttt{set\_option} do not occur. The kernel evaluation with \texttt{decide +kernel} occurs only in the data lemma \decl{a_data} of the statement file, which \decl{result} does not use.

\subsection*{Audit}
The file ends with the command
\begin{lstlisting}
#math0_settles Math0.Frontier.A392943.result Math0.Statements.A392943.statement
\end{lstlisting}
defined in \path{lean/Math0/Audit.lean} for statements that math0 formalizes itself. It fails unless (1) \decl{statement} is a definition of type \texttt{Prop} without universe parameters; (2) \decl{statement} and \decl{result} depend on no axioms other than \decl{propext}, \decl{Classical.choice} and \decl{Quot.sound} (this rejects \decl{sorryAx}, the axiom \decl{Lean.ofReduceBool} introduced by \texttt{native\_decide}, and custom axioms); and (3) the type of \decl{result} is syntactically the statement or its negation, after removing metadata annotations, with no unfolding. On this file it reports
\begin{lstlisting}
math0 audit: 'Math0.Statements.A392943.statement' depends on axioms [propext]
math0 audit: 'Math0.Frontier.A392943.result' depends on axioms [Classical.choice, Quot.sound, propext]
math0 statement check: 'Math0.Frontier.A392943.result' proves 'Math0.Statements.A392943.statement'
\end{lstlisting}

\subsection*{Toolchain and reproduction}
We use Lean~4 \cite{Lean4} \texttt{v4.33.1} and Mathlib \cite{Mathlib} \texttt{v4.33.1}, pinned in \path{lean/}; the repository also pins formal-conjectures \cite{FC} at commit \texttt{c77fdd4}, which this result does not use. To reproduce, run \path{scripts/verify.sh} from the repository root: it fetches the Mathlib build cache, builds all math0 results with \texttt{lake build} (re-running every audit command), checks that the pinned \texttt{formal\_conjectures} and \texttt{mathlib} checkouts are unmodified, and reports the number of passed statement checks. The single file can be checked with \texttt{lake env lean Math0/Frontier/OEIS/A392943.lean} in the directory \path{lean}; once the statement module is built, this needs well under a minute of CPU time, most of it spent loading the imported Mathlib modules (between about $6$ and $15$ seconds of user time in our runs). Building the statement module itself is dominated by the kernel check of \decl{a_data}, which takes well under a minute.

\begin{thebibliography}{99}

\bibitem{Aa}
M.~Aassila, \emph{Some results on Heron triangles}, Elem. Math. \textbf{56} (2001), no.~4, 143--146, \url{https://doi.org/10.1007/PL00000549}. Retracted by the editorial board of Elemente der Mathematik.

\bibitem{FC}
M.~Firsching, P.~Lezeau, S.~Mercuri, M.~Z. Horv\'ath, Y.~Dillies, C.~S\"onne, E.~Wieser, F.~Zhang, T.~Hubert, B.~Ag\"uera y Arcas, and P.~Kohli, \emph{Formal Conjectures: An open and evolving benchmark for verified discovery in mathematics}, arXiv:2605.13171, 2026. Repository: \url{https://github.com/google-deepmind/formal-conjectures}.

\bibitem{KL}
A.-V.~Kramer and F.~Luca, \emph{Some remarks on Heron triangles}, Acta Acad. Paedagog. Agriensis Sect. Mat. (N.S.) \textbf{27} (2000), 25--38.

\bibitem{Lean4}
L.~de Moura and S.~Ullrich, \emph{The Lean 4 theorem prover and programming language}, in: Automated Deduction -- CADE 28, Lecture Notes in Comput. Sci. \textbf{12699}, Springer, 2021, 625--635.

\bibitem{Mathlib}
The mathlib Community, \emph{The Lean mathematical library}, in: Proceedings of the 9th ACM SIGPLAN International Conference on Certified Programs and Proofs (CPP 2020), ACM, 2020, 367--381.

\bibitem{Read}
E.~Read, \emph{Integer triangles with integer circumradii}, Math. Gaz. \textbf{107} (2023), no.~569, 241--248, \url{https://doi.org/10.1017/mag.2023.55}.

\bibitem{vL}
R.~van Luijk, \emph{The diameter of the circumcircle of a Heron triangle}, Elem. Math. \textbf{63} (2008), no.~3, 118--121, \url{https://doi.org/10.4171/EM/96}.

\bibitem{A392942}
The OEIS Foundation Inc., \url{https://oeis.org/A392942}, entry by F.~Huber (2026).

\bibitem{A392943}
The OEIS Foundation Inc., \url{https://oeis.org/A392943}, entry by F.~Huber (2026).

\end{thebibliography}

\end{document}
